%% 第一章

\chapter{绪论}
\chaptermark{绪论}
\section{课题研究背景及意义}
近年来，全球范围内的太空探索活动不断推进，月球探测、火星巡视、小行星采样、空间站建设以及在轨航天器巡检等任务日益频繁，人类对深空地表环境和空间目标的观测需求持续增长。月球南极、火星岩石地表、小行星不规则表面等区域不仅具有重要科学价值，也直接关系到着陆区选择、巡视路径规划、采样点分析、任务仿真和设备安全评估等关键环节。随着深空探测任务从单纯的图像获取逐渐向环境理解、场景建模和任务辅助决策发展，如何从有限观测数据中恢复深空地表场景的空间结构，已经成为空间探测数据处理和智能感知中的重要研究问题。已有研究开始将传统摄影测量、神经辐射场和3D Gaussian Splatting等技术应用于月球、火星和小行星地表建模任务，说明深空地表三维重建与可视化正在成为相关领域的重要发展方向\upcite{ref6,ref7,ref8,ref9,ref10}。

在实际深空探测任务中，探测器与地面之间通常存在显著的通信延迟和带宽限制，图像数据的获取频率、分辨率和传输规模均受到任务条件约束。同时，探测器的观测位置和成像角度还受到轨道、姿态、能源、地形安全性以及任务规划等因素影响，使地面研究人员往往只能获得有限视角、有限分辨率和不同光照条件下的图像数据。二维图像虽然能够记录局部地表纹理、岩石分布和目标外观，但难以直接表达地形起伏、遮挡关系、空间尺度和复杂场景结构。当研究人员需要从多个角度观察目标区域、评估地形环境或进行任务仿真时，仅依赖二维图像往往难以形成直观、连续和可交互的空间认知。这就要求研究人员超越原始的视点局限，将零散的二维视觉信息升维转化为多维的空间表征。这种跨维度的信息整合在实际任务中发挥着核心枢纽作用：在着陆区选择与巡视路径规划中，三维场景能够为地形安全性评估、障碍物分析和路径可达性判断提供辅助信息；在采样点分析和科学解释中，三维重建结果有助于理解岩石形态、地貌结构和空间分布特征；在任务仿真与虚拟验证中，新视角合成能够支持从不同观察角度生成连续图像，为探测器操作规划、地面演示和环境可视化提供支撑。因此，围绕深空地表场景开展三维重建与新视角合成研究，对于提升深空探测数据的空间表达能力、环境理解能力和任务辅助决策能力具有重要工程应用价值。

与地面近景场景或室内场景相比，深空地表场景具有更加明显的非理想成像特征。深空探测图像中的场景往往具有极端深度跨度，同一幅图像中可能同时包含近处局部地形、远处地表结构以及尺度差异显著的空间目标，这使得场景几何关系和空间尺度表达更加复杂。同时，月面、火星表面和小行星表面常呈现低纹理、重复纹理或大面积颜色相近的区域，导致基于特征匹配的传统多视图几何方法容易出现匹配不足和结构恢复不完整等问题。此外，深空地表缺少稳定均匀的大气散射条件，光照方向、阴影范围、曝光状态和相机响应差异会造成明显的跨视角外观变化，进一步增加三维重建和新视角合成的难度。这些特点表明，深空地表场景的三维表达不能简单依赖面向常规地面场景的建模方法，而需要结合有限视角、弱纹理、复杂光照和大深度跨度等成像条件开展针对性研究。综上，研究深空地表场景的三维重建与新视角合成，不仅有助于提升深空探测数据的三维化表达和可视化分析能力，也能够为空间目标观察、远距离地表环境建模、航天场景仿真和虚拟视角生成等应用提供技术支撑。

\section{相关工作}
\subsection{从传统三维重建到NeRF与3DGS的发展}
传统三维重建主要建立在多视图几何基础之上。SfM方法通过特征提取、图像匹配、相机位姿估计、三角化和Bundle Adjustment恢复相机轨迹与稀疏点云，

代表性工作包括Photo Tourism和COLMAP等系统\upcite{ref1,ref2,ref3}。传统方法具有明确几何约束和较强可解释性，在纹理充足、视角覆盖充分、光照变化较小的场景中表现稳定。然而，当场景中存在低纹理区域、重复纹理、强阴影或视角覆盖不足等问题时，特征匹配和相机位姿估计的可靠性会受到影响，进而限制后续稠密重建和视觉合成效果。

NeRF将三维场景表示为连续的神经辐射场，通过体渲染方程从多视角图像中学习空间密度和颜色分布，显著推动了新视角合成质量提升\upcite{ref4}。与传统几何方法相比，NeRF能够通过隐式神经表示建模复杂的几何结构和视角相关外观，在新视角合成任务中表现出较强的表达能力。随后，相关研究也开始将神经场表示引入行星地表和深空场景建模中。例如，MaRF将NeRF思路引入火星地表表示，利用火星车和直升机图像合成火星环境，为行星地表建模提供了神经场方法的早期探索\upcite{ref7}。Asteroid-NeRF则针对小行星地表重建引入SDF和外观嵌入，以应对不同太阳高度角和曝光条件下的外观变化\upcite{ref8}。这些工作表明，神经表示能够补充传统摄影测量方法在低纹理、非均匀光照和复杂地表上的不足。

在神经渲染快速发展的基础上，3D Gaussian Splatting进一步提出使用显式三维高斯作为场景表示，在保持高质量渲染的同时获得实时渲染效率\upcite{ref5}。相比NeRF需要沿光线进行大量采样，3DGS通过屏幕空间splatting和alpha blending实现快速渲染，因而适合需要交互式观察的三维场景。随着3DGS的快速发展，围绕不同应用场景和技术瓶颈的改进工作不断涌现，例如面向大规模场景的CityGaussian、Scaffold-GS和Octree-GS，以及面向稀疏视角和位姿不足条件的InstantSplat、DUSt3R和VGGT；相关系统综述也对该领域的技术脉络、典型方法和未来趋势进行了归纳总结\upcite{ref6,ref11,ref12,ref13,ref14,ref15,ref16}。总体来看，从传统多视图几何到NeRF，再到3DGS，三维重建与新视角合成方法逐渐从显式几何恢复扩展到神经表示和显式可微渲染表示，为复杂场景下的高质量视觉建模提供了新的技术基础。

\subsection{非理想条件下三维重建的发展}
尽管NeRF和3DGS等方法显著提升了三维场景表示与新视角合成能力，但在真实复杂场景中，模型仍然会受到观测条件不足和成像条件变化的影响。对于本文关注的深空地表场景而言，数据采集往往存在视角有限、远近尺度差异明显、地表纹理弱、重复纹理多、阴影强烈以及曝光变化等问题。这些因素会削弱多视图之间的稳定约束，使模型在几何恢复、外观一致性和远区域重建等方面面临更大挑战。已有研究也开始比较摄影测量、NeRF与Gaussian Splatting在月球地形重建中的效果，表明神经表示和高斯表示能够在传统摄影测量基础上补充视觉细节和视角合成能力\upcite{ref9}。Martian World Models则进一步结合火星地表三维重建和可控视频合成，为火星任务仿真、机器人训练和导航规划提供了面向任务的三维世界模型思路\upcite{ref10}。这些研究说明，面向深空地表的三维重建已经逐渐从单纯几何恢复扩展到可视化、仿真和任务驱动建模，但其非理想观测条件仍然需要进一步处理。

稀疏视角和位姿不足是非理想三维重建中的典型问题。视角数量不足会导致多视图约束不充分，模型容易过拟合训练视角，并在新视角下产生漂浮物、几何破碎或纹理错位等问题。针对这一类问题，SCGaussian利用匹配先验增强少样本新视角合成中的结构一致性，InstantSplat通过密集立体先验和高斯优化缓解pose-free与sparse-view条件下的初始化问题，D2GS则从深度和密度角度分析稀疏视角下近处过拟合与远处欠拟合问题\upcite{ref14,ref17,ref24}。这些方法表明，在观测视角不足时，仅依赖图像重建损失往往难以获得稳定几何结构，引入深度、匹配或多视图先验是提升欠约束区域重建质量的重要途径。

外观变化也是非理想场景中影响三维重建质量的重要因素。在非受控图像集合或深空地表成像中，不同视角之间可能存在光照方向变化、曝光差异、阴影分布变化和局部亮度不一致等现象，这些因素会破坏传统多视图方法和3DGS优化中常用的光度一致性假设。针对这一问题，Gaussian in the Wild和WildGaussians分别通过外观分解、可见性建模和鲁棒特征处理非约束图像集合中的外观变化\upcite{ref18,ref19}。Luminance-GS使用视角自适应曲线调整处理低光、过曝和曝光变化，GaussHDR面向高动态范围场景学习统一的3D与2D局部色调映射\upcite{ref20,ref21}。Relightable 3D Gaussians则进一步将材质、BRDF和光照因素纳入高斯表示，用于更真实的重光照渲染\upcite{ref22}。这些研究说明，外观一致性已经成为3DGS在真实复杂场景中落地应用的关键问题。

除外观变化外，高斯形态稳定性也是3DGS在非理想条件下需要重点解决的问题。由于3DGS采用大量可优化的显式高斯基元表示场景，当视角约束不足或区域纹理信息较弱时，高斯分布可能出现不合理的形态变化，进而影响深度结构、表面连续性和新视角渲染质量。Depth-Consistent 3DGS通过物理散焦建模和多视图几何监督增强深度一致性，MVGS利用多视角联合训练和跨视角densification改善单视角训练带来的几何过拟合\upcite{ref23,ref25}。Effective Rank Analysis从高斯基元形态和几何退化角度分析了3DGS表示中的稳定性问题，为后续提升高斯几何质量提供了重要参考\upcite{ref26}。此外，3D-HGS从高斯基元形态角度探索更适合表面表示的半高斯结构\upcite{ref27}。这些工作表明，面向非理想条件下的3DGS重建，仅提升图像渲染质量是不够的，还需要从深度约束、外观一致性和高斯形态稳定性等方面共同改善场景表示质量。

\section{本文主要工作内容}
第一，搭建基于3D Gaussian Splatting的深空地表场景三维重建与新视角合成系统，实现从数据输入、模型训练、结果渲染到指标评估的完整流程。系统以多视角图像和相机参数为输入，通过显式三维高斯表示重建地表场景，并输出新视角渲染图像和三维高斯点云模型。

第二，设计FAME远区域增强模块。该模块利用辅助单目深度结果划分远区域，并在训练损失中对远区域重建误差进行额外约束，缓解常规图像重建损失容易偏向近处高纹理区域的问题。通过该设计，模型在优化过程中能够更加关注远处地表结构、目标边缘和背景层次。

第三，引入高斯基元形态稳定约束。本文根据高斯尺度参数设计有效秩正则与体积正则，用于约束高斯尺度分布和整体体积变化，降低高斯异常拉伸、过度膨胀或局部尺度突变对渲染质量和局部结构表达造成的不利影响。

第四，设计对光照扰动情况下的外观建模实验流程。针对深空地表场景可能出现的低光、过曝、白平衡变化和局部照明变化，本文结合外观分解分支分析复杂光照条件下的新视角合成效果。

第五，构建以新视角合成质量为核心的实验评价体系。本文围绕整体图像重建质量、远区域渲染效果、光照扰动条件下的外观一致性以及模块消融等方面设计实验，采用PSNR、SSIM和LPIPS对不同方法及本文各功能模块的性能表现进行定量评估。依托近六百个深空地表场景构成的大规模实验数据，本文开展了方法对比、远区域分析、光照扰动实验和消融实验，为验证所提方法的有效性和系统整体性能提供了实验依据。

\section{论文结构安排}
第一章为绪论。本章首先介绍深空地表场景三维重建与新视角合成的研究背景和研究意义，分析该任务在空间探测、地表环境建模、任务仿真和视觉感知等方面的应用价值。随后，围绕传统三维重建、NeRF与3D Gaussian Splatting等相关技术的发展脉络展开综述，并进一步分析深空地表场景在稀疏视角、远区域欠拟合、光照变化、弱纹理和几何不稳定等非理想条件下面临的主要挑战。

第二章为相关理论与技术基础。本章主要介绍本文研究所涉及的基础理论和关键技术，包括三维重建与新视角合成的基本概念、SfM与COLMAP的相机位姿估计流程、3D Gaussian Splatting的场景表示与可微渲染机制、单目深度估计在远区域辅助建模中的作用，以及PSNR、SSIM、LPIPS等新视角合成质量评价指标，为后续方法设计和系统实现提供理论基础。

第三章为系统需求分析与总体设计。本章从系统目标和应用需求出发，分析深空地表场景三维重建与新视角合成系统需要解决的核心问题，并对系统整体架构进行设计。具体包括数据输入与预处理、COLMAP位姿估计、3DGS场景初始化、模型训练、结果渲染、实验评估和结果统计等流程，同时对系统功能模块划分、运行方式和整体工作流程进行说明。

第四章为基于3DGS的深空地表场景重建方法。本章是全文的核心方法部分，重点介绍本文在3D Gaussian Splatting基础上进行的改进设计。首先说明基础3DGS重建流程，然后围绕深空地表场景中远区域欠拟合、几何结构不稳定和光照外观不一致等问题，分别介绍远区域划分策略、FAME远区域增强方法、几何正则约束以及光照外观建模方法，并说明各模块在整体训练流程中的作用。

第五章为实验设计与结果分析。本章围绕本文提出的方法和系统实现开展实验验证，介绍实验数据、实验环境、对比方法、消融设置和评价指标。在实验分析中，从整体图像重建质量、远区域表现、光照扰动条件下的外观一致性和模块贡献等角度对实验结果进行定量和定性分析。并通过对比实验与消融实验验证本文方法在深空地表场景三维重建与新视角合成任务中的有效性。

第六章为总结与展望。本章对全文所完成的研究工作进行总结，概括本文在深空地表场景三维重建与新视角合成方面的主要工作。同时，针对当前方法仍存在的不足，分析后续可以进一步改进和完善的方向，并对深空地表场景三维重建与新视角合成技术的未来发展作出展望。
